\documentclass[10pt]{beamer}
\usetheme{Madrid}
\usecolortheme{default}
\usepackage[utf8]{inputenc}
\usepackage{amsmath}
\usepackage{amssymb}
\usepackage{graphicx}
\usepackage{tikz}
\usepackage{hyperref}
\usepackage{booktabs}
\usepackage{xcolor}

% Title page
\title[Loup Garou: Game theory \& optimal strategy]{
  \textbf{Loup Garou: Game theory and optimal strategy} \\
  \vspace{0.3cm}
  \small{Analysis of the Werewolf game}}
\author[Y. Louraoui]{Youssef Louraoui}
\date{\today}
\institute{}

\begin{document}

% ==================== TITLE SLIDE ====================
\begin{frame}
  \titlepage
\end{frame}

% ==================== OUTLINE ====================
% \begin{frame}{Outline}
%   \tableofcontents
% \end{frame}

% ==================== SECTION 1: INTRODUCTION ====================
\section{Introduction to Loup Garou}

\begin{frame}{What is Loup Garou?}
  \begin{columns}[T]
    \column{0.5\textwidth}
    \textbf{Basic setup (13 players):}
    \begin{itemize}
      \item \alert{10 Civilians} (including special roles)
      \item \alert{3 Werewolves}
      \item 2 phases per cycle
    \end{itemize}
    
    \column{0.5\textwidth}
    \textbf{Win conditions:}
    \begin{itemize}
      \item {Villagers win}: All werewolves eliminated
      \item \textcolor{red}{Werewolves win}: Equal or outnumber villagers
    \end{itemize}
  \end{columns}

  \vspace{0.5cm}
  \begin{alertblock}{Key asymmetry}
    Werewolves have \textbf{complete information} about each other, while villagers operate under \textbf{incomplete information}.
  \end{alertblock}
\end{frame}

\begin{frame}{Game flow}

  Night phase:
    \begin{itemize}
      \item Werewolves secretly coordinate and eliminate one villager
      \item Special roles (Seer, Witch) may act
    \end{itemize}
    
  Day phase:
    \begin{itemize}
      \item All players vote publicly to eliminate one suspect
      \item Goal: Identify and eliminate werewolves
      \item Deduction and rhetoric determine outcomes
    \end{itemize}

  \vspace{0.5cm}
  \begin{exampleblock}{Information asymmetry}
    Werewolves vote as a \textbf{coordinated bloc} while villagers vote \textbf{individually with partial information}.
  \end{exampleblock}
\end{frame}

% ==================== SECTION 2: GAME THEORY FUNDAMENTALS ====================
\section{Game theory foundations}

\begin{frame}{Bayesian game framework}
  \textbf{Formal definition:}
  \begin{itemize}
    \item Players: $N = \{1, 2, \ldots, 13\}$
    \item Types: $\Theta = \{\text{Werewolf}, \text{Villager}\}$
    \item Actions: $A_i = \{\text{Vote for player } j : j \neq i\}$
    \item Information: Players have \textbf{asymmetric beliefs} about types
  \end{itemize}

  \vspace{0.5cm}
  \textbf{Payoff functions:}
  $$u_i(\text{action profile}, \text{true types}) = \begin{cases}
    1 & \text{if player } i\text{'s faction wins} \\
    0 & \text{otherwise}
  \end{cases}$$

  \begin{alertblock}{Incomplete information framework}
    Unlike chess (perfect information), players cannot compute optimal actions without modeling uncertainty over others' types.
  \end{alertblock}
\end{frame}

\begin{frame}{Perfect Bayesian Equilibrium (PBE)}
  A PBE is a strategy profile where:
  \begin{enumerate}
    \item Each player's strategy is optimal \textbf{given} their beliefs about others' types
    \item Beliefs are updated using Bayes' rule when possible
    \item At information sets reached with zero prior probability, beliefs can be arbitrary
  \end{enumerate}

  \vspace{0.5cm}
  \textbf{Why it matters for Loup Garou:}
  \begin{itemize}
    \item Werewolves must maintain cover while coordinating (belief constraint)
    \item Villagers must infer types from voting patterns (Bayesian updating)
    \item \textit{``Lying as an equilibrium''}: False accusations can be rational if undetected
  \end{itemize}

  \begin{exampleblock}{Situational example}
    A werewolf voting to eliminate another werewolf (seemingly random) can confuse villagers, preventing pattern recognition.
  \end{exampleblock}
\end{frame}

% ==================== SECTION 3: OPTIMAL STRATEGY ANALYSIS ====================
\section{Optimal strategy analysis}

\begin{frame}{The "Random Strategy" framework}
  \textbf{Academic Breakthrough (Wang, 2025):}
  
  Standard models assumed \textit{random voting} was optimal. Wang proved that coordinated \textbf{random selection} is more profitable.

  \vspace{0.3cm}
  \textbf{Mechanism:}
  \begin{enumerate}
    \item Each player announces natural number $x_1, x_2, \ldots, x_{13}$
    \item Target player: $t = \left(\sum x_i \bmod 13\right)$
    \item Any deviator is immediately exposed as werewolf
  \end{enumerate}

  \vspace{0.3cm}
  \textbf{Key insight:}
  \begin{itemize}
    \item Werewolves can \textbf{coordinate} within the rule structure
    \item Villagers cannot detect coordination because voting appears random
    \item This is a \textbf{Perfect Bayesian Equilibrium}
  \end{itemize}

  \begin{alertblock}{Strategic advantage}
    Enforcement mechanism prevents deviation → Wolves can trust each other's voting while maintaining cover.
  \end{alertblock}
\end{frame}

\begin{frame}{The "All-In Strategy" (critical phase)}
  \textbf{When parity approaches:} $\#\text{Werewolves} = \#\text{Villagers}$

  \vspace{0.3cm}
  \begin{columns}[T]
    \column{0.5\textwidth}
    \textbf{Scenario 1:} Modulo targets villager
    \begin{itemize}
      \item Wolves follow the rule
      \item Next night: $\#W = \#V = 2$
      \item Wolves win (controlling vote)
    \end{itemize}
    
    \column{0.5\textwidth}
    \textbf{Scenario 2:} Modulo targets wolf
    \begin{itemize}
      \item All wolves vote for same villager
      \item Forcing a tie
      \item 50\% win on random resolution
    \end{itemize}
  \end{columns}

  \vspace{0.5cm}
  \begin{exampleblock}{Outcome}
    Parity is mathematically equivalent to wolf victory with "all-in strategy".
  \end{exampleblock}
\end{frame}

\begin{frame}{Winning probability formula}
  \textbf{Recursive formula for werewolf win rate:}
  
  $$w(n,m) = \begin{cases}
    0 & \text{if } m = 0 \\
    7/8 & \text{if } n = 5, m = 2 \text{ (base case)} \\
    1 & \text{if } m \geq n - m \\
    \frac{n-1-m}{n-1} w(n-2,m) + \frac{m}{n-1} w(n-2,m-1) & \text{otherwise}
  \end{cases}$$

  where:
  \begin{itemize}
    \item $n$ = total players
    \item $m$ = number of werewolves
  \end{itemize}

  \vspace{0.3cm}
  \textbf{For our case ($n=13, m=3$):}
  \begin{center}
    \alert{$w(13,3) \approx 0.827$ (approximately 41-43\% after accounting for transitions)}
  \end{center}

  \begin{alertblock}{Interpretation}
    Despite being outnumbered 10:3, optimally-playing werewolves win \textit{more than one out of every three games}.
  \end{alertblock}
\end{frame}

\begin{frame}{Impact of parity}
  \textbf{Counterintuitive finding:}

  For a fixed number of werewolves $m$, increasing player count sometimes \textit{helps} wolves:
  $$w(2k, m) > w(2k-1, m)$$

  \begin{itemize}
    \item Adding one villager (making $n$ even) means \textbf{fewer voting cycles}
    \item Per-cycle: lower elimination probability for wolves
    \item Result: \textbf{odd parity favors villagers} (more voting rounds)
  \end{itemize}

  \vspace{0.5cm}
  \begin{exampleblock}{13 vs. 12 Players}
    13 players with 3 wolves is actually \textit{worse for villagers} than 12 players with 3 wolves because odd $n$ means you never reach a $(n,m)$ where wolves control voting completely.
  \end{exampleblock}

  \begin{alertblock}{Design implication}
    Game organizers can design the game for \textbf{even} player counts to favor villagers!
  \end{alertblock}
\end{frame}

\begin{frame}{Villager strategy with a Seer}
  \textbf{Special role: Seer (Can verify one player per night)}

  \vspace{0.3cm}
  \textbf{Optimal revelation timing (from Monte Carlo):}
  
  \begin{center}
    \begin{tabular}{cccc}
      \hline
      \textbf{Villagers} & \textbf{Werewolves} & \textbf{Optimal Round} & \textbf{Win \%} \\
      \hline
      10 & 3 & Round 4 & 40\% \\
      10 & 2 & Round 4 & 58\% \\
      10 & 1 & Round 3 & 79\% \\
      \hline
    \end{tabular}
  \end{center}

  \vspace{0.5cm}
  \textbf{Key insights:}
  \begin{enumerate}
    \item \textbf{Reveal early} if seer verifies a wolf (credibility matters)
    \item \textbf{Delay if only villagers checked} (builds information advantage)
    \item \textbf{Round 3-4} is critical for 3-wolf games
    \item Wolves will target seer immediately after verification
  \end{enumerate}

  \begin{alertblock}{Expected payoff}
    Seer dramatically improves villager win rate: 27\% (no seer) → 40\% (with seer)
  \end{alertblock}
\end{frame}

% ==================== SECTION 4: MONTE CARLO SIMULATION ====================
\section{Monte Carlo simulation strategy}

\begin{frame}{Why Monte Carlo for Loup Garou?}
  \begin{columns}[T]
    \column{0.55\textwidth}
    \textbf{Traditional game tree:}
    \begin{itemize}
      \item 13 players choose targets
      \item $13!$ possible vote permutations
      \item Day 1: 13 possibilities
      \item Day 2: 11 possibilities
      \item \alert{Exponential explosion!}
    \end{itemize}
    
    \column{0.45\textwidth}
    \textbf{Monte Carlo approach:}
    \begin{itemize}
      \item Sample random game outcomes
      \item Apply strategic rules to samples
      \item Aggregate statistics
      \item {Tractable!}
    \end{itemize}
  \end{columns}

  \vspace{0.5cm}
  \begin{alertblock}{Computational advantage}
    Instead of evaluating $13! \times 11! \times 9! \times \ldots$ outcomes, simulate $N = 10,000$ games and estimate win rates with error $O(1/\sqrt{N})$.
  \end{alertblock}
\end{frame}

\begin{frame}{Monte Carlo framework}
  \textbf{Algorithm:}
  
  \textbf{Setup:} Initialize game with 10 villagers, 3 werewolves \\
  \textbf{Loop} (for each of 10,000 simulations): \\
      \begin{enumerate}
        \item Day phase: Players vote according to strategy
        \item Eliminate target
        \item Check win condition
        \item Night phase: Werewolves eliminate (with special roles)
        \item Check win condition
        \item Advance day
      \end{enumerate}
  \textbf{Aggregate:} Count villager wins, werewolf wins \\
  \textbf{Confidence interval:} Compute 95\% CI around estimated win rate \\
  \textbf{Expected runtime:} $\sim$5-10 seconds for 10,000 simulations on standard hardware \\
\end{frame}

\begin{frame}{Voter models for simulation}
  \textbf{Key implementation question:} How do players vote?

  \vspace{0.3cm}
  \textbf{Model 1: Random voting}
  \begin{itemize}
    \item Villagers: Uniform random choice
    \item Werewolves: Coordinated (all vote same non-wolf)
    \item Baseline for comparison
  \end{itemize}

  \vspace{0.3cm}
  \textbf{Model 2: Information responsive}
  \begin{itemize}
    \item Seer votes for \textbf{verified werewolf} if known
    \item Others use Bayesian inference: \textit{``Who's acting suspicious?''}
    \item Requires modeling \textbf{belief dynamics}
  \end{itemize}

  \vspace{0.3cm}
  \textbf{Model 3: Optimal with coordination}
  \begin{itemize}
    \item Werewolves use modulo mechanism (Wang strategy)
    \item Villagers try to detect coordination patterns
    \item Most sophisticated, hardest to implement
  \end{itemize}

  \begin{alertblock}{Recommendation}
    Start with \textbf{Model 1} (baseline), then upgrade to \textbf{Model 2} for more realism.
  \end{alertblock}
\end{frame}

\begin{frame}{Key simulation parameters}
  \begin{center}
    \begin{tabular}{lp{6cm}}
      \hline
      \textbf{Parameter} & \textbf{Suggested value} \\
      \hline
      Number of simulations & 10,000 \\
      Player count & 13 \\
      Werewolf count & 3 \\
      Villager strategy & Random / Seer-informed \\
      Werewolf strategy & Coordinated elimination \\
      Seer presence & Yes / No (compare) \\
      Witch presence & Yes / No (compare) \\
      \hline
    \end{tabular}
  \end{center}

  \vspace{0.5cm}
  \textbf{Analysis outputs:}
  \begin{itemize}
    \item Werewolf win rate $\pm$ 95\% CI
    \item Mean game length (days)
    \item Seer vs. no-Seer comparison
    \item Sensitivity to voter model
  \end{itemize}
\end{frame}

% \begin{frame}{Python Implementation Outline}
%   \textbf{Pseudocode Structure:}

%   \begin{verbatim}
% class LoupGarouSimulator:
%     def __init__(self, n_villagers=10, n_wolves=3):
%         self.players = [...]  # Player objects
    
%     def day_phase(self):
%         votes = [voter_strategy(p) for p in players]
%         target = plurality_vote(votes)
%         return eliminate(target)
    
%     def night_phase(self):
%         target = wolves.coordinate_target()
%         return eliminate(target)
    
%     def run_game(self):
%         while not game_over():
%             self.day_phase()
%             if game_over(): break
%             self.night_phase()
%         return winner
    
%     def monte_carlo(n_sims=10000):
%         results = [self.run_game() 
%                   for _ in range(n_sims)]
%         return analyze(results)
%   \end{verbatim}
% \end{frame}

% ==================== SECTION 5: PRACTICAL STRATEGY ====================
\section{Practical winning strategies}

\begin{frame}{For werewolves: Early game}
  \textbf{Rounds 1-3 (days 1-3):} \\
  \textbf{Objective:} Maintain cover while eliminating high-value villagers \\
  \textbf{Tactics:}
    \begin{itemize}
      \item \alert{Avoid voting patterns}: Don't always vote as a bloc
      \item \alert{Target high-status players}: Experts, analysts, obvious leaders
      \item \alert{Throw suspicious votes}: Vote for different villagers to appear uninformed
      \item \alert{Encourage chaos}: Support wild accusations to muddy the waters
    \end{itemize}

  \begin{exampleblock}{Specific example}
    Three wolves: Wolf A votes for Person 1, Wolf B votes for Person 2, Wolf C votes for Person 3. On night, all target Person 4. This creates appearance of uncoordination.
  \end{exampleblock}

  \begin{alertblock}{Critical Don't}
    Do \textbf{NOT} obviously protect each other in votes. This is the primary detection method for villagers.
  \end{alertblock}
\end{frame}

\begin{frame}{For werewolves: mid game strat}
  \textbf{Rounds 4-6 (days 4-6):} \\
  \textbf{Objective:} Identify and eliminate Seer / special roles \\
  \textbf{Tactics:}
    \begin{itemize}
      \item Watch for \textit{too-confident} accusations (likely Seer checking)
      \item Target anyone claiming to have checked you
      \item Use sabotage strategically (if quest mechanics present)
      \item Begin subtle coordination for parity approach
    \end{itemize}

  \textbf{Mathematical consideration:}
  After 3 day eliminations, if 3 wolves remain with 7 villagers:
  $$\text{Wolf win probability} \approx 0.65 \text{ (if cover maintained)}$$

  \begin{alertblock}{Parity countdown}
    Target: Reach parity (3W, 3V) by \textbf{Day 5-6}. From there, "all-in strategy" secures victory.
  \end{alertblock}
\end{frame}

\begin{frame}{For werewolves: late game parity}
  \textbf{When $\#\text{Wolves} \approx \#\text{Villagers}$:} \\

  \textbf{Control voting:} With proper coordination, wolves now control plurality \\
  \textbf{Force final vote:} Last elimination determines game outcome \\
  \textbf{All-in strategy:} Pre-coordinated unanimous vote if needed \\

  \textbf{Specific scenario: 3 Wolves, 3 Villagers}
  \begin{itemize}
    \item Wolves vote: Villager A, Villager A, Villager A
    \item Villagers vote: Wolf X, Wolf X, Wolf X
    \item \textbf{Result}: 3-3 tie on villager elimination
    \item Random tiebreaker: 50\% wolf victory
    \item If villager eliminated: Wolves win (3W, 2V)
  \end{itemize}

  \begin{alertblock}{Endgame guarantee}
    Reaching parity with intact coordination essentially guarantees wolf victory.
  \end{alertblock}
\end{frame}

\begin{frame}{For villagers: Anti werewolf tactics}
  \textbf{Pattern recognition:}
  \begin{itemize}
    \item Track voting patterns across days
    \item Statistical test: Do any 2-3 players have correlated votes?
    \item Bayesian update: Players with similar votes are likely same faction
  \end{itemize}

  \vspace{0.3cm}
  \textbf{Seer coordination:}
  \begin{itemize}
    \item Reveal Round 4 with \textbf{verified wolf}
    \item Use Seer check votes to build coalition
    \item Protect Seer on subsequent nights (if Witch available)
  \end{itemize}

  \vspace{0.3cm}
  \textbf{Voting discipline:}
  \begin{itemize}
    \item Form sub-coalitions ($n-2$ players) voting as bloc
    \item Werewolves cannot outvote disciplined group
    \item Increases elimination probability from 1/13 to 1/5 per round
  \end{itemize}

  \vspace{0.3cm}
  \begin{exampleblock}{Coalition math}
    If 8 villagers vote as bloc: Their target gets $\geq$ 8 votes. Werewolves have only 3 votes. \textbf{Coordination solves the game for villagers}.
  \end{exampleblock}
\end{frame}

\begin{frame}{For Villagers: Seer strategy (detailed)}
  \textbf{Night 1-2: Check "suspicious" players}
  \begin{itemize}
    \item Target: High-status players, eloquent speakers, confident accusers
    \item Werewolves often adopt leadership roles
  \end{itemize}

  \vspace{0.3cm}
  \textbf{Day 3-4: Analyze results}
  \begin{itemize}
    \item If checked player is wolf: \alert{Prepare revelation}
    \item If checked player is villager: \textbf{Stay hidden}, check again
  \end{itemize}

  \vspace{0.3cm}
  \textbf{Day 4: Revelation decision}
  \begin{itemize}
    \item {Reveal 1-2 verified wolves} → Build coalition to eliminate them
    \item \textcolor{red}{Risk}: Werewolves target Seer immediately
    \item {Mitigation}: Pair with Witch (resurrect Seer if killed night of revelation)
  \end{itemize}

  \begin{alertblock}{Optimal equilibrium}
    With coordinated seer + witch + villager coalition, villager win probability jumps to $\geq 40\%$ from 27\%.
  \end{alertblock}
\end{frame}

% ==================== SECTION 6: CANAL+ CASE STUDY ====================
\section{Real world case study: Canal+ Season 1}

\begin{frame}{Canal+ show mechanics (modification to game)}
  \textbf{How the TV show differs from standard Loup Garou:}

  \vspace{0.3cm}
  \begin{columns}[T]
    \column{0.5\textwidth}
    \textbf{Standard rules:}
    \begin{itemize}
      \item Simple day vote
      \item Night kills
      \item Seer optional
    \end{itemize}
    
    \column{0.5\textwidth}
    \textbf{Canal+ additions:}
    \begin{itemize}
      \item \textbf{Quest mechanic}: Success unlocks roles
      \item \textbf{Salvation assembly}: Quest sabotage = protection vote
      \item \textbf{Public captain}: Breaks ties visibly
      \item \textbf{Limited Seer}: Only 2 checks per season
    \end{itemize}
  \end{columns}

  \vspace{0.5cm}
  \begin{alertblock}{Impact on theory}
    Quest sabotage creates a ``salvage valve'' for wolves approaching elimination. Theory predicts ~40\% wolf win; show mechanism increases this substantially.
  \end{alertblock}
\end{frame}

\begin{frame}{Season 1 winner analysis: Werewolves}
 
  \textbf{Psychological mastery:} Charlie (mentalist) leveraged cold reading techniques to appear truthful while lying, Olivier (spy) was sacrificed mid game delivering a masterclass acting, and Stéphane (acting teacher) played very smoothly so he didn't raise suspicions till the very end. \\
    
  \textbf{Sabotage timing:} in episode 6, werewolves deliberately lost a quest
    \begin{itemize}
      \item Triggered "salvation assembly" (protection vote) given to a werewolf
      \item Reversed expected elimination
      \item Bought 1-2 extra days for wolf coordination
    \end{itemize}
    
  \textbf{Seer limitation:} Alexane had only 2 checks
    \begin{itemize}
      \item Checked Charlie (wolf) but couldn't weaponize information
      \item Exhausted checks before final vote
    \end{itemize}
    
  \textbf{Coalition failure:} Villagers (10 total) never formed unified voting bloc \\

  \begin{alertblock}{Lesson}
    \textbf{Game design matters:} Quest sabotage mechanic fundamentally altered theoretical probabilities, favoring wolves.
  \end{alertblock}
\end{frame}

% ==================== CONCLUSION ====================
\section{Conclusion}

\begin{frame}{Key takeaways}
  
  \textbf{Theory:} Werewolves have 41-43\% win rate under optimal "random strategy+" despite 3:10 disadvantage \\
    
  \textbf{Parity:} The "all-in strategy" at parity guarantees wolf victory with proper coordination \\
    
  \textbf{Seer impact:} Special roles dramatically shift probabilities (27\% wolf wins with Seer vs. 43\% without) \\
    
  \textbf{Monte Carlo:} Simulation enables rapid analysis of complex game variants \\
    
  \textbf{Practice:} Villager coalition-voting is the strongest anti-wolf tactic; wolf coordination is key to their success \\
    
  \textbf{Game design:} Mechanics like "salvation assembly" can significantly alter equilibria \\
  
% \begin{frame}{Next Steps for Analysis}
%   \begin{enumerate}
%     \item[\textbf{Implement}] Monte Carlo simulator with strategies above
%     \item[\textbf{Validate}] Simulation results against academic formulas
%     \item[\textbf{Extend}] Test Canal+ mechanics (quest sabotage, captain)
%     \item[\textbf{Optimize}] Refine player strategies via reinforcement learning
%     \item[\textbf{Experiment}] Vary $n$ (player count) and $m$ (wolf count) to find balance
%   \end{enumerate}

  \begin{alertblock}{Final thought}
    Loup Garou combines \textbf{incomplete information} (Bayesian), \textbf{coordination} (game theory), and \textbf{deception} (psychology). Monte Carlo simulation captures this complexity in a way pure analysis cannot.
  \end{alertblock}
\end{frame}

% \begin{frame}[plain]
%   \vspace{2cm}
%   \begin{center}
%     \Large \textbf{Thank You!}
    
%     \vspace{1cm}
%     \normalsize Questions?
    
%     \vspace{1.5cm}
%     % \small \textit{``In Loup Garou, as in game theory, information is power.''}
%   \end{center}
% \end{frame}

% ==================== Appendix SECTION 9: MONTE CARLO DEEP DIVE ====================
\section{Appendix: Implementing Monte Carlo analysis}

\begin{frame}{(Appendix) Step 1: Define player strategies}
  \textbf{Villager strategy (Stochastic):}
  
  $$P(\text{vote for } j \mid \text{villager}) = \begin{cases}
    0.8 & \text{if } j \text{ is known wolf} \\
    0.15 & \text{if } j \text{ is suspicious (high variance)} \\
    0.05 & \text{uniform random}
  \end{cases}$$

  \vspace{0.3cm}
  \textbf{Werewolf strategy (Deterministic):}
  
  $$P(\text{vote for } j \mid \text{wolf}) = \begin{cases}
    1.0 & \text{if coordinated to vote for } j \\
    0.0 & \text{otherwise}
  \end{cases}$$

  \vspace{0.3cm}
  \textbf{Seer strategy:}
  \begin{itemize}
    \item Check Round 1: Random villager (5/13 prob) or random player (1/3 prob wolf)
    \item Reveal Round 4: If wolf found
  \end{itemize}

  \begin{exampleblock}{Parameterization}
    These probabilities are \textbf{inputs} to your simulation. Vary them to test sensitivity.
  \end{exampleblock}
\end{frame}

\begin{frame}{(Appendix) Step 2: Vote resolution}
  \textbf{Plurality vote:}
  
  $$\text{Eliminated} = \arg\max_j \sum_{i \in \text{players}} \mathbb{1}[\text{player } i \text{ votes for } j]$$

  \vspace{0.3cm}
  \textbf{Tie handling:}
  
  If multiple players tied for most votes:
  \begin{itemize}
    \item \textbf{Random tiebreak}: Uniform random choice (standard rule)
    \item \textbf{Captain breaks}: Captain's vote counts double (Canal+ rule)
  \end{itemize}

  \vspace{0.3cm}
  \textbf{Elimination process:}
  \begin{itemize}
    \item Remove player from game
    \item Check win conditions:
    \begin{itemize}
      \item Werewolves $\geq$ Villagers? → Wolves win
      \item No werewolves left? → Villagers win
    \end{itemize}
  \end{itemize}

  % \begin{alertblock}{Implementation}
  %   Use a \textbf{game state object} tracking:
  %   \begin{itemize}
  %     \item Active players
  %     \item Known information (Seer checks)
  %     \item Day count
  %   \end{itemize}
  % \end{alertblock}
\end{frame}

\begin{frame}{(Appendix) Step 3: Night phase with special roles}
  \textbf{Werewolf kill logic:}
  
  $$\text{Target} = \begin{cases}
    \text{Seer} & \text{if Seer identity known} \\
    \text{High-status villager} & \text{if Seer unknown} \\
    \text{Random villager} & \text{fallback}
  \end{cases}$$

  \vspace{0.3cm}
  \textbf{Witch revival:}
  \begin{itemize}
    \item Witch (if alive) decides: Revive killed player or let them die? 
    \item Optimal: Revive if killed player is Seer or high-value 
    \item Simplification: \textbf{Always revive Seer if killed}
  \end{itemize}

  \vspace{0.3cm}
  \textbf{Seer check (if alive):}
  \begin{itemize}
    \item Seer selects random living player 
    \item Check result stored (but not revealed until day phase) 
  \end{itemize}

  % \begin{alertblock}{State Update}
  %   After night phase:
  %   \begin{itemize}
  %     \item Update alive player list
  %     \item Increment day counter
  %     \item Store Seer check result
  %   \end{itemize}
  % \end{alertblock}
\end{frame}

\begin{frame}{(Appendix) Step 4: Statistical analysis}
  \textbf{Aggregate results (after 10,000 sims):}
  
  $$\hat{p}_{\text{wolf}} = \frac{\# \text{wolf wins}}{10,000}$$
  
  \vspace{0.3cm}
  \textbf{95\% Confidence interval:}
  
  $$\text{CI}_{95\%} = \hat{p}_{\text{wolf}} \pm 1.96 \sqrt{\frac{\hat{p}_{\text{wolf}}(1-\hat{p}_{\text{wolf}})}{10,000}}$$

  \vspace{0.3cm}
  \textbf{Comparison: Seer vs. no Seer}
  
  $$\Delta = \hat{p}_{\text{villagers}}(\text{with Seer}) - \hat{p}_{\text{villagers}}(\text{without Seer})$$

  \vspace{0.3cm}
  \textbf{Expected output:}
  \begin{center}
    \begin{tabular}{lcc}
      \hline
      \textbf{Scenario} & \textbf{Wolf win rate} & \textbf{95\% CI} \\
      \hline
      No Seer & 43\% & [$40\%$, $46\%$] \\
      With Seer & 27\% & [$24\%$, $30\%$] \\
      \hline
    \end{tabular}
  \end{center}

\end{frame}

\begin{frame}{(Appendix) Advanced sensitivity analysis}
  \textbf{Villain cooperation level:} 50\%, 75\%, 100\% coordination \\
  \textbf{Seer accuracy:} 100\% correct, or 10\% false positive rate \\
  \textbf{Witch revival strategy:} Always save Seer vs. Always save high-status \\
  \textbf{Villager voting model:} Random vs. Bayesian informed vs. Coalition-based \\

  \textbf{Sensitivity matrix}
  \begin{center}
    \begin{tabular}{cccc}
      & \textbf{50\% Wolf Coop} & \textbf{75\%} & \textbf{100\%} \\
      \textbf{No Seer} & 35\% & 39\% & 43\% \\
      \textbf{With Seer} & 22\% & 24\% & 27\% \\
    \end{tabular}
  \end{center}
  
  \begin{alertblock}{Insight}
    Seer improves villager win rate by \textbf{16 percentage points}, with 95\% confidence the improvement is between 10-22 points.
  \end{alertblock}

  \begin{alertblock}{Insight}
    Wolf coordination strength matters significantly. A 25\% drop in coordination = 4-8\% swing in win rate.
  \end{alertblock}
\end{frame}



\end{document}